% One-page general ML / computer vision résumé (website: public/files/Michael_Perez_Resume.pdf).
% Every fact comes from documents/cv/cv.tex, src/data/content.js, or the owner's
% job-agent experience bank (~/Desktop/job-agent/bank). No proficiency is inferred.
%
% Layout: the type size, leading, and gaps below are tuned so the content fills
% exactly one letter page. After editing, rebuild and check the page count; the
% knobs to re-fit are listed in README.md.
\documentclass[10pt,letterpaper]{article}
\usepackage{fontspec}
\setmainfont{Times New Roman}
\usepackage[left=0.7in,right=0.7in,top=0.55in,bottom=0.55in]{geometry}
\usepackage{enumitem}
\usepackage[hidelinks,unicode]{hyperref}
\hypersetup{pdftitle={Michael F. Pérez — ML and Computer Vision Engineer},pdfauthor={Michael F. Pérez}}
\pagestyle{empty}
\setlength{\parindent}{0pt}
\setlength{\parskip}{0pt}
\setlength{\emergencystretch}{2em}
\raggedright
\hyphenpenalty=10000 % never split a word across lines: keeps copy-paste and ATS text extraction clean
\setlist[itemize]{leftmargin=0.17in,labelsep=0.07in,itemsep=1.5pt,parsep=0pt,topsep=2.5pt,partopsep=0pt}

\newcommand{\ressection}[1]{\par\addvspace{9pt}\textbf{\MakeUppercase{#1}}\par\nobreak\vspace{1.5pt}\hrule height 0.5pt\nobreak\vspace{4.5pt}}
\newcommand{\resrow}[2]{\noindent #1\hfill\mbox{#2}\par}
% \resrole{title}{employer, place}{dates}
\newcommand{\resrole}[3]{\par\addvspace{5pt}\resrow{\textbf{#1}, #2}{#3}}
% \resskill{label}{items}
\newcommand{\resskill}[2]{\textbf{#1:} #2\par\vspace{1pt}}
% \respub{title}{role, venue, link}
\newcommand{\respub}[2]{\textbf{#1.} #2\par\vspace{2.5pt}}

\begin{document}
\fontsize{10.5}{12.9}\selectfont
\begin{center}
{\fontsize{19}{22}\selectfont\textbf{MICHAEL F. PÉREZ}}\\[3pt]
\textbf{ML \& Computer Vision Engineer} · Ph.D. in Computer Science · Tampa, FL (open to relocation)\\[1.5pt]
\href{mailto:michaelperez023@gmail.com}{michaelperez023@gmail.com} | (954) 882-5069 | \href{https://www.michaelfperez.com}{michaelfperez.com}\\[1.5pt]
\href{https://www.linkedin.com/in/michaelperez023/}{linkedin.com/in/michaelperez023} | \href{https://github.com/michaelperez023}{github.com/michaelperez023} | \href{https://scholar.google.com/citations?user=fCzyXhUAAAAJ\&hl=en}{Google Scholar}
\end{center}
\vspace{-4pt}
ML and computer vision engineer (Ph.D. in Computer Science, University of Florida, 2026) with experience integrating computer vision into production software, building ML data and evaluation workflows, and researching video understanding across two DARPA-funded programs.

\ressection{Technical Skills}
\resskill{ML \& computer vision}{PyTorch, TensorFlow, OpenCV, scikit-learn, Hugging Face, VLMs (LaViLa, LLaVA, Qwen)}
\resskill{MLOps \& data tools}{Kubernetes, Airflow, MLflow, FiftyOne, Label Studio, Docker, SLURM/HPC}
\resskill{Languages \& systems}{Python, C++, SQL, JavaScript, FastAPI, Flask, Celery, Redis, Linux, Git, NumPy, pandas, SciPy}

\ressection{Experience}
\resrole{Senior Engineer}{Raft · Tampa, FL}{Aug 2026 – Nov 2026}
\begin{itemize}
\item Developed and integrated computer vision data exploration and analysis capabilities into a larger production software system, including dataset visualization, clustering, filtering, natural-language querying, and geospatial analytics.
\item Built ML data management and evaluation workflows for dataset curation, annotation review, model analysis, and debugging.
\item Contributed to the development, testing, and integration of AI/ML capabilities, including object-detection experimentation and improvements to existing software systems.
\end{itemize}
\resrole{Machine Learning Engineer Intern}{CoVar · Durham, NC}{May 2025 – Aug 2025}
\begin{itemize}
\item Applied infrared video super-resolution to long-range object classification; first-author paper at the 2026 Military Sensing Symposium (MSS).
\item Trained lightweight RGB object-detection models targeting real-time, on-device drone detection.
\item Built offline computer vision evaluation pipelines spanning preprocessing, inference, metrics, and efficiency analysis.
\end{itemize}
\resrole{Graduate Research Assistant}{University of Florida · Gainesville, FL}{Aug 2020 – Aug 2026}
\begin{itemize}
\item Built UI and image/video analysis APIs and trained action-segmentation models for CReLeRI, a DARPA-funded explainable long-video analysis system; first author of its ACM Multimedia 2025 paper.
\item Built CReLeRI's long-video analysis web tool: FastAPI backend with Celery and Redis background tasks, containerized with Docker, with GPU compute offloaded to UF's HiPerGator cluster.
\item Trained and tested action-recognition models, including fine-tuning the video-language model LaViLa, for ENKIx, a DARPA-funded real-time AR task-guidance system, and ran usability and cognitive-load evaluations.
\item Ph.D. thesis: developed efficient action-detection models for long videos and built a unified framework benchmarking detectors across mAP, latency, throughput, memory, and FLOPs; first-author manuscript in preparation.
\end{itemize}

\ressection{Education}
\resrow{\textbf{Ph.D. in Computer Science}, University of Florida}{Aug 2026}
\resrow{\textbf{M.S. in Computer Science}, University of Florida}{May 2023}
\resrow{\textbf{B.S. in Computer Science}, Florida Polytechnic University}{May 2020}

\ressection{Selected Publications}
\respub{CReLeRI: Explainable, Concept-centric, Representation, Learning, Reasoning, and Interaction Video Analysis System}{First author · ACM Multimedia 2025 · \href{https://doi.org/10.1145/3746027.3754479}{doi:10.1145/3746027.3754479}}
\respub{Super-Resolution for Improving Vehicle Recognition at Large Stand-offs}{First author · MSS 2026}
\respub{MuCHEx: A Multimodal Conversational Debugging Tool for Interactive Visual Exploration of Hierarchical Object Classification}{Co-author · IEEE Computer Graphics and Applications, 2025 · \href{https://doi.org/10.1109/MCG.2025.3598204}{doi:10.1109/MCG.2025.3598204}}

\ressection{Honors \& Service}
\textbf{Honors:} L3Harris Graduate Fellowship (one of five UF CISE recipients, 2022); UF Graduate School Preeminence Award\par\vspace{1pt}
\textbf{Service:} Peer reviewer for ACM Multimedia, ACM ICMR, IEEE VR, and ICXR (2024 – 2026)
\end{document}
